\begin{table}[H]
\centering
\caption{Convergencia beta de la remuneración laboral departamental, 2010--2025}
\label{tab:dept_remuneracion_convergencia_24}
\small
\begin{tabular}{lrrrrl}
\toprule
Indicador & $\beta$ & Error est. & Estad. $t$ & $R^2$ & Lectura \\
\midrule
Remuneración por trabajador & 0,0089 & 0,0073 & 1,22 & 0,064 & No concluyente \\
Remuneración por hora & 0,0060 & 0,0082 & 0,73 & 0,024 & No concluyente \\
\bottomrule
\end{tabular}
\caption*{\footnotesize Nota: la variable dependiente es el crecimiento anualizado del indicador y la variable explicativa es el logaritmo del nivel inicial del mismo indicador. Un coeficiente $\beta<0$ indica convergencia beta; un coeficiente $\beta>0$ indica que los departamentos con mayor remuneración inicial crecieron más rápido. La lectura exige que el coeficiente sea estadísticamente distinto de cero al 5\%. Fuente: cálculos propios con GEIH.}
\end{table}